%% Working-Paper-Vorlage für das ML4SCS-Semesterprojekt.
\documentclass[11pt]{article}
\PassOptionsToPackage{ngerman}{babel} %% Deutsche Silbentrennung; die Vorlage lädt babel selbst.
\usepackage{UF_FRED_paper_style}
%% Die Vorlage laedt palatino (Type1-Familie ppl), die unter der
%% Unicode-Engine von tectonic nicht existiert: LaTeX ersetzt sie durch die
%% Standardschrift und verliert dabei die fette Serie (\textbf, \paragraph).
%% TeX Gyre Pagella ist der freie Palatino-Nachbau mit allen Schnitten.
\usepackage{fontspec}
\setmainfont{texgyrepagella}[
    Extension    = .otf,
    UprightFont  = *-regular,
    BoldFont     = *-bold,
    ItalicFont   = *-italic,
    BoldItalicFont = *-bolditalic]

\usepackage{tikz}
\usetikzlibrary{arrows.meta, positioning, calc}

\onehalfspacing
\setlength{\droptitle}{-5em} %% Vorgabe der bereitgestellten Vorlage.

\title{Erkennung von Schreibaktivität anhand der IMU-Daten einer Apple Watch}

\author{Noah Samel (4004067)\\
    noah.samel@stud.leuphana.de
\and
    Tajuddin Snasni (TODO: Matrikelnummer)\\
    tajuddin.snasni@stud.leuphana.de}

\date{2. September 2026}

\begin{document}
\selectlanguage{ngerman}

{\setstretch{.8}
\maketitle
\begin{abstract}
\noindent
Schreiben ist eine alltägliche Tätigkeit. Diese Arbeit
untersucht, ob sich Schreibaktivität allein aus den Daten der IMU - Daten
einer Smartwatch am Schreibhandgelenk erkennen lässt,
unabhängig von der tragenden Person. Dafür wurde ein eigenes Erfassungssystem
aus Apple Watch, iPhone-Bridge und Server entwickelt. Ein Moleskine Smart Pen
lieferte die Ground Truth, deren Strichereignisse über eine varianzbasierte
Zeitausrichtung mit den Watch-Daten verknüpft werden. Unter einem
Latin-Square-balancierten Studienprotokoll mit Schreibaufgaben und gezielt
schreibähnlichen Gegenaufgaben (Tippen, Scrollen, Gestikulieren) entstand ein
Datensatz mit 32 Personen und rund 11,5 Stunden annotierter Daten. Alle
Modelle wurden mit einer Five-Fold Cross-Validation über die
Personen bewertet, sodass keine Person je in ihrem eigenen Training war. Ein Random
Forest auf 88 Merkmalen je Sekunde erreicht 0,865 Accuracy und 0,942
ROC-AUC. Ein nachgeschaltetes Zwei-Zustands-HMM hebt dieselben Vorhersagen ohne
Neutraining auf 0,917 Accuracy und 0,973 ROC-AUC, wenn es die Aufnahme
rückblickend dekodiert, und auf 0,898 als kausaler Filter für den
Live-Einsatz. Ein zeitliches Faltungsnetz mit rekurrentem Kopf (TCN-BiGru)auf
5-s-Fenstern erreicht 0,921 Accuracy und 0,975 ROC-AUC. Wir schließen,
dass ein personenunabhängiger Schreibdetektor auf einer Smartwatch möglich
ist und dass weitere Fortschritte eher von vielfältigeren Trainingspersonen
als von neuen Modellen zu erwarten sind.

\medskip
\noindent
\textit{\textbf{Schlagwörter: }Aktivitätserkennung; Inertialsensorik;
Wearable Computing; Schreiberkennung; maschinelles Lernen.}
\end{abstract}
}

% Tajis Abschnitte
\section{Einleitung}
\label{sec:introduction}

Wearable-Geräte sind aus dem Alltag nicht mehr wegzudenken. Smartwatches werden durchgehend am Handgelenk getragen und zeichnen kontinuierlich Bewegungsdaten über integrierte Inertialmesseinheiten (IMUs) auf, bestehend aus Beschleunigungssensoren und Gyroskopen. Diese Sensoren wurden ursprünglich für Schrittzählung und Fitness-Tracking entwickelt, erfassen jedoch ein weit breiteres Spektrum an Aktivitätsinformationen. Eine bislang kaum untersuchte Anwendung ist die automatische Erkennung von Schreibaktivität: die Frage, ob eine Person gerade mit der Hand schreibt, allein auf Basis der Handgelenksbewegung.

Handschreiben ist eine grundlegende Tätigkeit in Bildung, Beruf und Alltag. Im Vergleich zum Tippen auf einer Tastatur ist Handschreiben eine kognitiv anspruchsvollere motorische Aufgabe, die das Gehirn auf andere Weise beansprucht und nachweislich die Informationsverarbeitung fördert. Die Fähigkeit, Schreibaktivität passiv und unauffällig zu erkennen, ohne dass der Nutzer aktiv einen Timer starten oder mit dem Gerät interagieren muss, eröffnet vielfältige praktische Anwendungen. Dazu zählen die Erfassung produktiver Schreibzeit, das Monitoring kognitiver Aktivität in Bildungskontexten sowie der Aufbau smarter Umgebungen, die sich daran orientieren, ob eine Person gerade gedanklich etwas zu Papier bringt.

Bisherige Arbeiten zur Schreiberkennung mit handgelenksgetragenen Sensoren fokussieren sich vorwiegend darauf, \textit{was} geschrieben wird, also auf die Erkennung von Buchstaben, Wörtern oder Schreibgesten, anstatt der einfacheren, aber praktisch relevanteren Frage nachzugehen, \textit{ob} überhaupt geschrieben wird \cite{soubam2023, xiamotionhacker2018}. Diese Ansätze basieren häufig auf kontrollierten Laborbedingungen, wenigen Probanden und Evaluationsstrategien, die keine belastbare Aussage über die Generalisierbarkeit auf unbekannte Nutzer ermöglichen.

Die vorliegende Arbeit adressiert diese Einschränkungen. Wir erheben IMU-Daten einer Apple Watch Series~7 von 22 Probanden in einem ausbalancierten Studienprotokoll, das sowohl echte Schreibbedingungen als auch eine Reihe schwieriger Negativbeispiele umfasst, darunter Tastatur- und Smartphone-Tippen, Scrollen, Stiftfummeln und Gestikulieren. Die Ground-Truth-Labels werden ausschließlich während der Datenerhebung von einem Moleskine Smart Pen geliefert; nach dem Training läuft die Inferenz allein auf der Apple Watch. Alle Modelle werden mit Leave-One-Subject-Out-Kreuzvalidierung (LOSO) bewertet, bei der in jeder Falte ein Proband vollständig ausgelassen wird. Dies entspricht exakt der Generalisierungsanforderung eines realen Einsatzszenarios.

Unsere zentrale Forschungsfrage lautet:

\begin{quote}
\textit{Lässt sich Schreibaktivität allein aus den IMU-Daten einer Apple Watch zuverlässig erkennen, unabhängig von der schreibenden Person und dem geschriebenen Inhalt?}
\end{quote}

Wir vergleichen verschiedene Ansätze, darunter einen Random-Forest-Basisklassifikator mit 92 handkonstruierten Features, einen kausalen Hidden-Markov-Modell-Nachfilter (HMM) sowie tiefe Zeitreihenmodelle (TCN, GRU). Die experimentellen Ergebnisse und eine detaillierte Analyse der Modellleistung werden in Abschnitt~4 vorgestellt.

Der Rest des Artikels gliedert sich wie folgt: Abschnitt~2 bespricht verwandte Arbeiten zur Schreiberkennung mit handgelenksgetragenen Sensoren. Abschnitt~3 beschreibt das Datenerhebungsprotokoll, den Sensor-Setup und die Vorverarbeitungs-Pipeline. Abschnitt~4 präsentiert die experimentellen Ergebnisse und deren Analyse. Abschnitt~5 schließt mit einer Diskussion der Limitationen und zukünftiger Forschungsrichtungen.

\section{Verwandte Arbeiten}
\label{sec:related_work}

Im Rahmen dieser Arbeit haben wir zwei Studien identifiziert, die unserem Projekt am nächsten kommen. Soubam et al. (2023) verwenden dieselbe Sensorklasse und denselben Klassifikationsansatz zur Erkennung von Schreibaktivität am Handgelenk. Xia et al. (2018) nutzen ebenfalls eine handgelenksgetragene IMU und einen Random Forest, verfolgen jedoch ein grundlegend anderes Ziel. Im Folgenden werden beide Arbeiten vorgestellt und in Bezug zu unserem Ansatz gesetzt.

\subsection{Erkennung von Schreib-Mikroereignissen (Soubam et al., 2023)}

Soubam et al.~\cite{soubam2023} präsentieren ein System zur Erkennung feinkörniger Schreib-Mikroereignisse mithilfe einer Smartwatch. Anstatt den geschriebenen Inhalt zu erkennen, untersuchen sie, welche Teilaktivität eine Person gerade ausführt. Die drei betrachteten Klassen sind aktives Schreiben (\textit{write}), Verschieben der Hand auf dem Papier (\textit{shift}) sowie Beginn einer neuen Zeile (\textit{newline}). Diese Formulierung ist methodisch nah an unserer binären Schreiberkennungsaufgabe, operiert jedoch auf einer feineren Granularitätsstufe.

Die Daten wurden von zehn Probanden (fünf Frauen, fünf Männer) mit einer LG W100 Smartwatch erhoben, die Beschleunigungssensor-, Gyroskop- und Rotationsvektor-Daten bei 100~Hz aufzeichnet. Jede Testperson absolvierte drei Schreibszenarien, nämlich Abschreiben, Schreiben aus dem Gedächtnis und freies Formulieren, über drei Sitzungen. Dies ergab insgesamt rund 15 Stunden annotierter Sensordaten. Die Ground-Truth-Labels wurden durch manuelle Videoannotation mit der ELAN-Software erstellt.

Ein zentraler methodischer Beitrag des Papers ist das \textit{Thresholding Tag Assignment}-Verfahren zur Fensterbeschriftung. Da sich Mikroereignisse stark in ihrer Dauer unterscheiden, wobei \textit{write}-Segmente mehrere Sekunden dauern, während \textit{shift} und \textit{newline} oft unter einer Sekunde liegen, klassifiziert die übliche Mehrheitsabstimmung kurze Ereignisse systematisch falsch. Das vorgeschlagene Verfahren weist einem Fenster stattdessen das Label jedes Ereignisses zu, das einen definierten Mindestanteil von beispielsweise 20\% überschreitet, wobei seltenere Klassen Vorrang haben. Dadurch können mittlere Fenstergrößen von 200 bis 300 Samples verwendet werden, ohne kurze Ereignisse in der Labelverteilung zu verlieren.

Zur Klassifikation werden handkonstruierte Feature-basierte Modelle (Random Forest, SVM, XGBoost, MLP) mit einem LSTM auf Rohzeitreihen verglichen. Statistische Features wie Mittelwert, Standardabweichung, Kurtosis, Zero-Crossing-Rate und Signal-Magnitude-Area sowie Frequenz-Domain-Features wie FFT-Energie und Schiefe werden für alle Sensorachsen extrahiert. Die Evaluation erfolgt mittels fünffacher Kreuzvalidierung mit dem Macro-F1-Score. Das beste nutzerunabhängige Modell erreicht F1~=~0,84, ein nutzerspezifisches Modell F1~=~0,91. Bemerkenswert ist, dass der Random Forest konsistent mit dem LSTM gleichzieht oder dieses übertrifft. Dies zeigt, dass sorgfältiges Feature Engineering bei kleinen Sensordatensätzen mit Deep Learning konkurrieren kann.

Unser 88-Feature-Set folgt unabhängig davon einer ähnlichen Logik und kombiniert Zeitbereichs-Statistiken mit FFT-basierten Spektral-Features. Unser Ansatz unterscheidet sich jedoch in drei wesentlichen Punkten. Erstens formulieren wir die Aufgabe als binäre Klassifikation (Schreiben vs.\ Nicht-Schreiben), was die Labelstruktur vereinfacht, aber anspruchsvolle Negativbeispiele wie Tastatur- und Smartphone-Tippen einführt. Zweitens erheben wir Daten von 32 Probanden, mehr als dreimal so viele wie Soubam et al. Drittens evaluieren wir mit einer nach Personen gruppierten Five-Fold Cross-Validation statt Standard-k-Fold über Fenster, sodass keine Person je in ihrem eigenen Training liegt. Das liefert eine deutlich konservativere und realistischere Schätzung der Generalisierung auf unbekannte Nutzer.

\subsection{MotionHacker: Belauschen von Handschrift (Xia et al., 2018)}

Xia et al.~\cite{xiamotionhacker2018} nähern sich dem Thema Handschreiben und Bewegungssensoren aus einer völlig anderen Perspektive, nämlich aus dem Bereich Datenschutz und Sicherheit. Ihr System MotionHacker ist ein Proof-of-Concept-Angriff, der zeigt, dass eine bösartige App auf einer Smartwatch unbemerkt Handgelenksbewegungen mitloggen und daraus rekonstruieren kann, was die Person schreibt, konkret welche Buchstaben und Wörter beim normalen Handschreiben auf Papier entstehen.

Das System zielt auf nutzerunabhängige Buchstabenerkennung für Kleinbuchstaben in Druckschrift ab. Eine LG R Watch zeichnet Beschleunigungssensor- und Gyroskop-Daten bei 200~Hz auf. Fünf Trainingspersonen schreiben alle 26 Buchstaben je 20-mal; fünf weitere Testpersonen fungieren in der Evaluation als Vergleichsgruppe. Die Pipeline gliedert sich in vier Stufen: Vorverarbeitung mit Ausreißerentfernung mittels 3$\sigma$-Kriterium und 1--25~Hz-Tiefpassfilter, Segmentierung der Wort- und Buchstabengrenzen über Gyroskop-Amplitudenschwellenwerte, Feature-Extraktion und Klassifikation.

Das Feature Engineering ist besonders aufwändig. Da die Handgelenksbewegung beim Schreiben subtil ist und stark von Schreibgeschwindigkeit und Druck abhängt, werden klassische Zeitbereichs-Features als unzureichend betrachtet. Stattdessen konstruieren die Autoren 115 spezialisierte, ausschließlich gyroskopbasierte Features in zwei Gruppen: 78 Frequenzbereichs-Features auf Basis von FFT-Entropie und Spektrum im 1--25~Hz-Band sowie 37 phasenbezogene Features, die die zeitlichen Beziehungen zwischen Peaks und Tälern über die drei Gyroskopachsen kodieren. Ein Random-Forest-Klassifikator liefert die Top-5-Buchstabenkandidaten je Segment; ein Wörterbuch-Filter ordnet anschließend Buchstabensequenzen dem wahrscheinlichsten Wort zu.

Evaluiert an den fünf Testpersonen, die zwei 100-Wörter-Absätze abschreiben, erreicht MotionHacker eine durchschnittliche Wort-Accuracy von 32,8\% für Top-1-Vorhersagen und 48,8\% für Top-5. Diese Ergebnisse bestätigen, dass IMU-Daten am Handgelenk ausreichend Signal tragen, um feinkörnige Schreibmuster nutzerunabhängig zu unterscheiden.

Unser Ansatz teilt mit MotionHacker die Sensormodalität (handgelenksgetragene IMU, Beschleunigungssensor und Gyroskop), den Kernklassifikator (Random Forest) und die Nutzung von FFT-basierten Frequenzbereichs-Features. Unser Ziel ist jedoch grundlegend verschieden und einfacher: Wir fragen nur, \textit{ob} geschrieben wird, nicht \textit{was}. Diese Vereinfachung erlaubt robustere Features und eine Evaluation auf einer größeren Kohorte.

\subsection{Einordnung unserer Arbeit}

Beide Studien belegen, dass die IMU am Handgelenk genug Signal für Schreibmuster trägt, und beide arbeiten wie wir mit handkonstruierten Zeit- und Frequenzbereichs-Features und einem Random Forest. Soubam et al. zeigen zudem, dass ein Random Forest auf Sensordaten kleiner Kohorten mit Deep Learning mithalten kann, ein Befund, den unsere Ergebnisse auf 1-s-Fenstern bestätigen und auf 5-s-Fenstern relativieren (Abschnitt~\ref{sec:results}). Gleichzeitig teilen beide Arbeiten dieselbe Schwäche: Beide evaluieren zwar personenunabhängig, aber auf zehn beziehungsweise fünf Testpersonen, zu wenig für eine belastbare Aussage über die Generalisierung auf unbekannte Nutzer, und keine von beiden prüft schreibähnliche Alltagsbewegungen als Negativklasse. Unsere Arbeit schließt diese Lücke durch eine Kohorte von 32 Probanden, ein ausbalanciertes Studiendesign mit gezielt schreibähnlichen Negativbeispielen und eine nach Personen gruppierte Five-Fold Cross-Validation, die der Bedingung eines realen Deployments auf der Uhr eines neuen Nutzers direkt entspricht.


% Noahs Abschnitte
\section{Data Collection}
\label{sec:data-collection}

% OWNER: Noah

\subsection{Sensing System}
% Apple Watch: accelerometer, gyroscope, capture rate and channel variants.
% Moleskine Smart Pen: used only to obtain ground-truth stroke events.
% State explicitly that inference requires only the watch.

\subsection{Study Protocol}
% Participants, tasks, counterbalancing, writing and non-writing activities.
% Mention Study Mode and task markers only where methodologically relevant.

\subsection{Dataset}
% Final cohort size, sessions, pool differences, inclusion/exclusion criteria,
% and class balance. Insert a compact cohort table once values are frozen.

\section{Methodology}
\label{sec:methodology}

% OWNER: Noah

\subsection{Clock Alignment and Ground-Truth Labels}
% Describe pen--watch offset recovery, watch-base merge, tolerance, and the
% capture-clock correction. Avoid implementation details that do not affect validity.

\subsection{Preprocessing and Windowing}
% Stable timestamp ordering, gap closing, 1 s windows and 0.5 s stride.

\subsection{Feature-Based Classifier}
% Summarise the 88 feature set, per-session normalisation, and Random Forest.

\subsection{Deep-Learning Models}
% Describe only architectures retained for the final comparison.

\subsection{Temporal Decoding and Model Fusion}
% Distinguish causal online filtering from non-causal daily smoothing.
% Include cross-paradigm fusion only if it appears in the final results.

\section{Versuchsaufbau und Evaluation}
\label{sec:experimental-setup}

\subsection{Evaluationsprotokoll}
\label{sec:cv}

Alle Kennzahlen stammen aus einer gruppierten Five-Fold Cross-Validation über
die Personen, abgestimmt mit der Partnergruppe an der ZHAW. Sämtliche Fenster
einer Person liegen entweder im Training oder im Test, jede der 32 Personen
ist genau einmal Testperson. Random Forest und Netze wurden mit
unterschiedlichen Fold-Partitionen gerechnet; gepaarte Tests zwischen den
Familien sind auf dieser Kohorte daher nicht möglich. Ein früher verwendetes
Leave-One-Subject-Out lieferte auf denselben Daten praktisch identische
Werte.

Glättung, Kalibrierung und Fehleranalyse rechnen auf der Out-of-Fold-Tabelle
der Fensterwahrscheinlichkeiten und erben dieselbe Leakage-Garantie; die
Übergangsmatrix des HMM wird je Runde nur aus den Trainingspersonen
geschätzt.

\subsection{Metriken und statistische Analyse}
\label{sec:metrics}

Primäre Kennzahl ist die Accuracy je Fenster bei fester Schwelle 0,5, weil
die Klassen nahezu ausgeglichen sind und ein Tracker eine Entscheidung
treffen muss. Ergänzend berichten wir ROC-AUC und den F1-Wert der Klasse
Schreiben, jeweils mit ihrer Entscheidungsskala.

Die Streuung zwischen Personen liegt bei drei bis vier Prozentpunkten.
Jeden Vergleich auf denselben Personen prüfen wir deshalb mit dem
Wilcoxon-Vorzeichen-Rang-Test \citep{wilcoxon1945} auf den personenweisen
Differenzen und nennen, bei wie vielen Personen eine Konfiguration vorn
liegt. Ein Gewinn unter einem Prozentpunkt ohne $p < 0{,}05$ gilt als
Rauschen. Für die Netze mitteln wir drei Seeds, weil zwei GPU-Läufe mit
identischem Seed um bis zu 1,6 Prozentpunkte auseinanderliegen können; bei
32 Personen liegen für die Netze nur Seed-Aggregate vor, ihr Vergleich
untereinander ist dort nominal.

\paragraph{Modellauswahl auf derselben Kohorte.} Architekturen und
Hyperparameter wurden anhand derselben Cross-Validation ausgewählt, auf der
wir berichten. Personen kreuzen nie zwischen Training und Test, aber die Werte der Netze
sind Selektionssieger und damit nach oben verzerrt; eine Nested
Cross-Validation war mit dem Rechenbudget nicht darstellbar. Die Sieger der
Hyperparametersuche verloren beim Nachrechnen mit drei Seeds ein bis zwei
Prozentpunkte gegenüber ihrem besten Einzellauf. Die Hyperparameter des
Random Forest blieben über das Projekt unverändert.

\subsection{Baselines und Ablationsstudien}
\label{sec:ablations}

Der Random Forest auf 88 Merkmalen ist die Referenz. Die Netze auf nativen
5-s-Fenstern vergleichen wir mit ihm bei gleicher Entscheidungslatenz: gegen
seine kausal über 5~s geglätteten 1-s-Vorhersagen und gegen einen Random
Forest mit direkt über 5~s berechneten Merkmalen. Die Ablationen ändern je
eine Variable auf identischen Folds und Seeds: Schwerkraftvektor,
Gyroskop, Dekodiermodus und Fusion. Zwei Vergleiche sind nicht auf eine
Variable beschränkt (50 gegen 100~Hz, Drei-Kanal-Lauf der Netze) und werden
entsprechend eingeordnet. Wirkungslose Eingriffe fasst Abschnitt~\ref{sec:ablation-results} zusammen,
weil sie die Interpretation der Leistungsgrenze tragen. Random Forest und HMM
laufen lokal in unter einer Stunde, die Netze brauchten auf einer
Cloud-GPU rund zehn Stunden.

\section{Results}
\label{sec:results}

% OWNER: Noah

\subsection{Writing-Activity Classification}
% Lead with the final grouped-five-fold headline result.
% TODO: Add the canonical model-comparison table generated from final artefacts.

\subsection{Temporal Decoding}
% Report causal filtering and non-causal smoothing separately.

\subsection{Ablation Results}
% Gravity, deep--deep fusion, HMM hyperparameters, and other retained analyses.

\subsection{Error Analysis}
% Keyboard/phone tapping, task-specific errors, and soft-writer case.
% Prefer participant- and task-level evidence over anecdotal descriptions.

\section{Diskussion}
\label{sec:discussion}

Die Forschungsfrage lautete, ob sich Schreibaktivität allein aus den
IMU-Daten einer Smartwatch erkennen lässt, unabhängig von der Person,
die sie trägt. Die Antwort ist ja, mit einer klar benennbaren Grenze. Auf 32
Personen, von denen jede für ihr eigenes Testergebnis nie im Training war,
erreicht ein zeitliches Faltungsnetz auf 5-s-Fenstern 92\,\% Accuracy, und
ein Random Forest auf handgebauten Merkmalen mit kausalem HMM-Filter 90\,\%.
Beide Zahlen sind ohne Kenntnis der Person, ohne Kalibrierungsphase und ohne
Blick in die Zukunft erreichbar. Was fehlt, sind die restlichen acht bis zehn
Prozent, und die Ergebnisse sagen ziemlich genau, wo sie sitzen.

\paragraph{Die Grenze des Merkmalsmodells liegt vermutlich in der
Repräsentation.} Über den größten Teil des Projekts sah es so aus, als läge
die Leistungsgrenze im Signal. Random Forest, MiniRocket, ein vortrainiertes
Modell und die Netze auf 1-s-Fenstern trafen dieselbe Marke und scheiterten
an denselben Personen. Kein Merkmalseingriff, keine Augmentation und keine
Umgewichtung bewegte sie. Der Schritt von 22 auf 32 Personen trennt diese
Lesart, wenn auch nicht sauber: Der Random Forest blieb stehen, die Netze
stiegen — zugleich änderten sich Kohortenzusammensetzung und Protokollanteil,
sodass die Beobachtung nicht allein der Personenzahl zuzuschreiben ist. Die
naheliegende Deutung ist, dass nicht das Signal erschöpft ist, sondern das,
was Statistiken über ein 1-s-Fenster ausdrücken können. Ein Netz, das die
Reihenfolge der Bewegungen innerhalb von fünf Sekunden sieht, profitiert
anders als der Random Forest weiter von zusätzlichen Personen. Dass die Netze im Modern-Pool auf
23 Personen den Schwerkraftvektor nutzen können, während der Random Forest
auf vier Personen daran scheiterte, passt in dasselbe Bild.

\paragraph{Zeitkontext ist der entscheidende Hebel, und man kann ihn nur
einmal einsammeln.} Drei ganz verschiedene Wege führen auf dieselbe Höhe:
ein Netz auf nativen 5-s-Fenstern, ein 1-s-Modell mit HMM, und die Fusion
von Random Forest und Netz. Das HMM hilft nur Modellen ohne eigenes
Gedächtnis. Sobald das Basismodell fünf Sekunden Kontext trägt, überglättet
es. Dass der Rolling-Mean auf keiner Skala hilft, das HMM aber auf jeder,
zeigt, dass es nicht um Mittelung geht, sondern um ein Modell der
Übergänge: Schreiben ist ein Zustand mit Trägheit, und ein Detektor, der das
weiß, darf bei schwacher Evidenz beim letzten Zustand bleiben. Für den
Einsatz ist damit die Frage nach der Latenz entscheidend. Wer eine Anzeige
im Moment will, bekommt 0,90. Wer abends den Tag rekonstruiert, bekommt
0,92, weil der Smoother beide Seiten eines Übergangs sehen darf. Diese
Differenz von knapp zwei Prozentpunkten ist in 31 von 32 Personen dieselbe
und war bei 22 Personen schon vorhanden. Sie gehört zu den robustesten
Befunden der Arbeit.

\paragraph{Der Restfehler ist personen- und aufgabenspezifisch, nicht
diffus.} Die Fehler des Modells liegen nicht gleichmäßig über allen
Nicht-Schreib-Situationen, sondern konzentrieren sich auf zwei Muster: Tippen
auf Tastatur oder Handy bei einem Teil der Personen und sehr leises oder
stark unterbrochenes Schreiben bei einem anderen Teil. Beide Muster sind am
Handgelenk physikalisch begründet. Wer beim Tippen die Hand ähnlich bewegt
wie beim Schreiben, erzeugt ähnliche Rotationen. Wer sehr leicht schreibt,
erzeugt kaum mehr Bewegung als in Ruhe. Für das Tippen konnten wir
zeigen, dass die trennende Information in den Merkmalen vorhanden ist, aber
von der Mehrheit der Merkmale überstimmt wird, und dass alle Versuche, sie
per Merkmalsdesign oder Gewichtung durchzusetzen, gescheitert sind. Der
gepoolte Tipp-Fehler ist mit den zusätzlichen Personen von 36 auf 22
Prozent gefallen, ohne dass jemand daran gearbeitet hätte. Mehr Personen mit
verschiedenen Tippstilen sind offenbar der Hebel, den Merkmalsdesign nicht
ersetzen kann.

\paragraph{Was die Zahlen messen, entscheidet die Labeldefinition.} Das
Schließen von Lücken bis 2,5~s macht aus der Frage \emph{berührt der Stift
das Papier} die Frage \emph{ist die Person im Schreibmodus}. Für einen
Schreibzeit-Tracker ist das die richtige Frage, und sie ist für das Modell
lernbar, weil das Handgelenk in Mikropausen dasselbe tut wie in den Strichen.
Die Kehrseite ist sichtbar: Rechenaufgaben mit Denkpausen über 2,5~s
bleiben schwierig, weil die Semantik dort ausfranst, und ein Schätzer für den
Schreibanteil überschätzt die rohe Stiftzeit systematisch um rund 20
Prozentpunkte. Beides ist kein Modellfehler, sondern eine Folge der
Definition, und beides muss bei jeder Anwendung mitgedacht werden.

\paragraph{Wo die Ergebnisse tragen und wo sie Hypothesen bleiben.} Belegt
sind, jeweils gepaart über Personen und mit $p < 0{,}05$: der Gewinn des
HMM-Filters, der Vorsprung des Smoothers, der Vorsprung der nativen
5-s-Netze gegenüber dem Random Forest bei 15 und 20 Personen, der Gewinn der
Fusion, die Kosten des fehlenden Gyroskops und die Wirkung der Korrekturen
an der Zeitachse. Nominal, aber nicht belegt sind die Rangfolge der beiden
Netze, der Vorsprung der Netze bei 32 Personen (der Abstand ist groß, die
personenweisen Vorhersagen liegen aber nicht vor) und die Wirkung des
Schwerkraftvektors auf tcn6. Hypothesen sind die Erklärungen: dass der
rekurrente Kopf Lageinformation aus dem Verlauf rekonstruiert, und dass die
Tipp-Verwechslung mit mehr Personen weiter schrumpft. Wir haben im Verlauf
des Projekts zweimal einen scheinbar sauberen Befund zurücknehmen müssen,
einmal wegen eines Sortierfehlers, einmal wegen eines verschobenen
Zeitversatzes. Beide Male war nicht die Statistik falsch, sondern die Daten
darunter. Das ist der Grund, warum diese Arbeit so viel Platz auf
Zeitachsen und Kontrollen verwendet.

\section{Limitationen}
\label{sec:limitations}

\paragraph{Kohorte und Aufgaben.} 32 Personen tragen eine
personenunabhängige Aussage, aber die Kohorte ist homogen: überwiegend
Studierende, ein Raum, ein Uhrenmodell, die Uhr immer am Schreibhandgelenk.
Händigkeit, Alter und Schreibgewohnheit wurden nicht erfasst; über Kinder,
ältere Menschen oder Personen mit motorischen Einschränkungen sagen die
Ergebnisse nichts. Alle Daten stammen aus vorgegebenen Aufgaben von ein bis
vier Minuten. Im Alltag sind Schreibphasen seltener und weniger klar
begrenzt, und die Negativklasse ist breiter als sechs Nicht-Schreibaufgaben.
Ein Schreibanteil von 43\,\% ist eine Folge des Protokolls; bei schiefen
Anteilen schrumpft ein Schätzer des Schreibanteils zur Mitte. Eine
Ganztagsaufzeichnung mit anderer Ground Truth (Tagebuch, Video) steht aus.

\paragraph{Zwei Datenpools.} Alle Hauptergebnisse laufen auf dem Legacy-Pool,
in dem 23 der 32 Aufnahmen heruntergerechnet sind. Der Dezimierer ist
phasenneutral, und ein explorativer Vergleich ergab keinen Unterschied, aber
der Pool bleibt eine Mischung aus nativen und abgeleiteten Signalen. Die
Aussagen zum Schwerkraftvektor stützen sich auf die 23 Personen des
Modern-Pools und sind nicht mit den 32-Personen-Zahlen vergleichbar.

\paragraph{Ground Truth aus dem Stift.} Das Toleranzfenster von $\pm 40$~ms
und das Schließen von Lücken bis 2,5~s sind Festlegungen, und beide
erzeugen inhärent mehrdeutige Fenster. Zwei Aufnahmen trugen über Wochen
einen falschen Zeitversatz, bevor eine Plausibilitätsgrenze ihn aufdeckte.
Eine videobasierte Kontrolle der Labels fehlt.

\paragraph{Statistische Absicherung der Netze.} Die Läufe auf 32 Personen
haben keine personenweisen Vorhersagen gespeichert; Rangfolge der Netze und
Abstand zum Random Forest sind dort nicht gepaart getestet. Die
Nichtdeterminiertheit der GPU (bis 1,6 Prozentpunkte bei identischem Seed)
begrenzt, wie fein Architekturen unterscheidbar sind, und die Modellauswahl
fand auf der berichteten Kohorte statt (Abschnitt~\ref{sec:cv}).

\paragraph{Zielgröße und Einsatz.} Das Modell unterscheidet Schreiben von
allem anderen, nicht Handschrift von Skizze, und sagt nichts über
Aufmerksamkeit; der Schreibanteil je Aufgabe ist nur ein Proxy. Der Random
Forest mit HMM-Filter läuft im Live-Dashboard. Ob die Hintergrund-Sensorik
der Watch über zwölf Stunden 50~Hz liefert und das konvertierte Modell auf
der Uhr dieselben Vorhersagen erzeugt, ist noch nicht gemessen.

\section{Fazit und Ausblick}
\label{sec:conclusion}

Schreibaktivität lässt sich aus den IMU-Daten einer Apple Watch am
Schreibhandgelenk personenunabhängig erkennen. Mit 32 Personen, einem
instrumentierten Stift als Ground Truth und einer Evaluation ohne Person im
eigenen Training erreicht ein Faltungsnetz mit rekurrentem Kopf auf
5-s-Fenstern 0,921 Accuracy und 0,975 ROC-AUC (drei Seeds, ohne gepaarten
Test). Ein Random Forest auf 88 Merkmalen erreicht je Sekunde 0,865, mit
kausalem HMM 0,898 und rückblickend 0,917; das HMM läuft ohne Neutraining im
Live-System.

Drei Erkenntnisse gehen über die Zahlen hinaus. Erstens liegt die Grenze des
Merkmalsmodells eher in der Repräsentation als im Signal: Zwölf zusätzliche
Personen haben den Random Forest nicht bewegt, die Netze schon. Zweitens ist
Zeitkontext der wirksamste Hebel und lässt sich nur einmal einsammeln, mit
zwei Prozentpunkten Aufschlag für alles, was nicht live entschieden werden
muss. Drittens sitzen die verbleibenden Fehler bei bestimmten Personen in
bestimmten Aufgaben, vor allem beim Tippen, und dort ist Trainingsvielfalt
der Hebel, nicht Merkmalsdesign.

Der nächste Schritt liegt deshalb in den Daten: Personen aufnehmen, die anders
tippen und schreiben als die bisherige Kohorte, und einen Teil der Aufnahmen
gegen Video annotieren, damit die Labeldefinition selbst überprüfbar wird.
Beides verspricht mehr als weitere Rechenzeit.


\medskip
\bibliography{references.bib}

\end{document}
